\documentclass[11pt]{article}
\usepackage[T1]{fontenc}
\usepackage[utf8]{inputenc}
\usepackage{lmodern,amsmath,amssymb,amsthm,mathtools}
\usepackage[margin=1in]{geometry}
\usepackage{microtype,booktabs,longtable,array,enumitem,xcolor,listings}
\usepackage[hidelinks]{hyperref}
\hypersetup{pdftitle={Puiseux Expansions in Intuitionistic Primitive Recursive Arithmetic},
 pdfauthor={Research report prepared for Vladimir Reshetnikov}}
\setlist{itemsep=3pt,topsep=5pt}
\setlength{\emergencystretch}{3em}
\newtheorem{theorem}{Theorem}[section]
\newtheorem{lemma}[theorem]{Lemma}
\newtheorem{proposition}[theorem]{Proposition}
\newtheorem{corollary}[theorem]{Corollary}
\theoremstyle{definition}
\newtheorem{definition}[theorem]{Definition}
\newtheorem{remark}[theorem]{Remark}
\newcommand{\N}{\mathbb N}
\newcommand{\Z}{\mathbb Z}
\newcommand{\Q}{\mathbb Q}
\newcommand{\C}{\mathbb C}
\newcommand{\A}{\overline{\mathbb Q}}
\newcommand{\ipra}{\ensuremath{\mathsf{iPRA}}}
\newcommand{\ord}{\operatorname{ord}}
\newcommand{\disc}{\operatorname{Disc}}
\newcommand{\Res}{\operatorname{Res}}
\newcommand{\coeff}{\operatorname{coeff}}
\newcommand{\trunc}{\operatorname{trunc}}
\newcommand{\PR}{\operatorname{PR}}
\newcommand{\norm}[1]{\left\lVert #1\right\rVert}
\newcommand{\abs}[1]{\left|#1\right|}
\newcommand{\repo}{\texttt{ProveIt}}
\newcommand{\pin}{\texttt{d851229c33ef395db8dc04baf28c6a291073f61e}}
\pdfstringdefDisableCommands{\def\ipra{iPRA}}
\lstset{basicstyle=\ttfamily\small,breaklines=true,columns=fullflexible,
 frame=single,rulecolor=\color{black!20},showstringspaces=false}
\title{\textbf{Puiseux Expansions in Intuitionistic\\
Primitive Recursive Arithmetic}\\[6pt]
\large A complete bounded proof for exact algebraic data,\\
constructive convergence, and a Lean formalization blueprint}
\author{Prepared for Vladimir Reshetnikov}
\date{15 November 2026}
\begin{document}
\maketitle
\begin{abstract}
For a polynomial supplied by exact complex-algebraic coefficient codes and an
exact algebraic expansion point, we construct all Puiseux branches with their
multiplicities by primitive recursive operations whose correctness is proved
by quantifier-free induction. The output includes a common ramification
\((d!)^d\), finite negative parts, and rational convergence radii and tail bounds.
The proof uses a discriminant to bound every coefficient search in a
degree-decreasing Newton--Hensel construction. Only one finite root jet is
needed: it yields a polynomial simple-root equation, after which a coefficient
recurrence and finite majorants produce all coefficients and prove convergence.

The input qualification is essential. A uniform procedure taking arbitrary
primitive recursive Cauchy coefficient names and returning exact Puiseux
coefficient names would decide halting, already for a monic quadratic. We prove
this obstruction and state the correct first-order arithmetical result.
An explicit exact-algebraic backend is supplied; its one root-existence
dependency is the human-readable constructive FTA proof in \repo, frozen in
this archive. The Lean deliverable specifies a deeply embedded \ipra\ target
and the remaining implementation obligations. It is an uncompiled outline,
not a completed Lean proof.
\end{abstract}
\newpage
\tableofcontents
\newpage

\section{The theorem being proved, and its constructive boundary}
\label{sec:scope}
The linked MathWorld statement concerns \emph{convergent} fractional-power
expansions of local algebraic functions, including Laurent expansions for
branches tending to infinity \cite{mathworld}. A formal splitting theorem
alone does not establish that analytic conclusion.

First-order \ipra\ has numerical variables and primitive recursive function
symbols, intuitionistic first-order logic, and quantifier-free induction.
It does not quantify over all complex numbers or all infinite sequences.
A precise result must specify the finite input, the functions producing
coefficients, and the finite identities and estimates expressing the theorem.

\subsection{Exact input and the main result}
An \emph{exact algebraic code} represents a complex algebraic number by a
pair of real-algebraic codes. A real code consists of a rational polynomial,
a rational isolating interval, and finite rational certificates under which
bisection names its unique root. Rational values also have direct codes.
Section~\ref{sec:backend} constructs all operations and certificates.
Equality of codes means decidable equality of the represented values.

\begin{theorem}[Effective convergent Puiseux theorem in \ipra]
\label{thm:main}
Let \(P(Z,Y)\in\A[Z,Y]\), \(d=\deg_Y P>0\), and \(z_0\in\A\) be supplied
by exact algebraic codes. Write
\[
 P(z_0+X,Y)=\sum_{r=0}^d a_r(X)Y^r.
\]
Fixed primitive recursive procedures return:
\begin{enumerate}
\item \(N=(d!)^d\), a pole bound \(p\in\N\), distinct branch indices
      \(1\le i\le n\), and positive multiplicities \(m_i\) with
      \(\sum_i m_i=d\);
\item primitive recursive exact algebraic coefficient functions
      \(w_{i,k}\), with notation
      \[
      W_i(T)=\sum_{k\ge0}w_{i,k}T^k,\qquad Y_i(T)=T^{-p}W_i(T);
      \]
\item positive rational numbers \(\eta,M_i\) such that, for every finite \(L\),
      \begin{equation}
       \sum_{k=0}^L |w_{i,k}|\eta^k\le M_i;
       \label{eq:main-weight}
      \end{equation}
\item certificates for the coefficientwise identity
      \begin{equation}
       \sum_{r=0}^d T^{(d-r)p}a_r(T^N)V^r
       =a_d(T^N)\prod_{i=1}^n(V-W_i(T))^{m_i}.
       \label{eq:main-cleared}
      \end{equation}
\end{enumerate}
For \(|t|\le q\eta\), \(0\le q<1\), and \(K>L\), the finite tail estimate is
\begin{equation}
 \left|\sum_{k=L+1}^K w_{i,k}t^k\right|\le M_iq^{L+1}.
 \label{eq:main-tail}
\end{equation}
The series consequently converge with explicit moduli and, on a sufficiently
small punctured disc,
\[
 P(z_0+t^N,Y)=a_d(t^N)\prod_i(Y-Y_i(t))^{m_i}.
\]
The finite coefficient identities and rational versions of the displayed
estimates are theorems of ordinary first-order \ipra.
\end{theorem}

The coefficient functions are coherent: later precision requests do not
change earlier coefficients. All coefficients lie in a single finite
algebraic extension of the field generated by the input. Only finitely many
root adjunctions occur in the initial construction; later recurrences use
field arithmetic alone.

Singular points, repeated factors, and poles are included. A branch has a
finite limit exactly when its first \(p\) numerator coefficients vanish,
a finite decidable test. Its limit is \(w_{i,p}\), and
\[
 Z=z_0+t^N,\qquad Y=w_{i,p}+\sum_{k\ge1}w_{i,p+k}t^k.
\]
The formulation concerns branches over a punctured \(X\)-disc. A component
contained entirely in \(X=0\) is not a graph of an algebraic function of \(X\);
polynomial content in \(X\) does not create another punctured-disc branch.

\subsection{Why arbitrary Cauchy coefficient names do not suffice}
\label{subsec:obstruction}
\begin{proposition}[Uniform coefficient-name obstruction]
\label{prop:obstruction}
No uniform computable procedure takes arbitrary primitive recursive Cauchy
coefficient names of a monic polynomial and always returns a positive
ramification and computable names for the coefficients of even one exact
Puiseux root.
\end{proposition}
\begin{proof}
Consider \(P_a(X,Y)=Y^2-X^2-a\), \(a\ge0\). A returned root after \(X=T^e\)
satisfies \(y(T)^2=T^{2e}+a\).
If \(a=0\), finite coefficient induction gives \(y_j=0\) for \(j<e\) and
\(y_e^2=1\), so \(|y_e|=1\). Any finite negative Laurent part first vanishes
by the same coefficient induction.
If \(a>0\), the constant coefficient satisfies \(y_0^2=a\) and is apart
from zero. Coefficient induction gives \(y_j=0\) for \(1\le j<2e\), so
\(y_e=0\).
Approximating \(y_e\) within \(1/8\) distinguishes the cases by a rational
comparison of the squared modulus with \(1/4\).

For a machine index \(v\), let \(H(v,k)\) indicate that the machine halts
by step \(k\). The one fixed primitive recursive family
\[
 A(v,s)=\sum_{k<s}2^{-k-1}H(v,k)
\]
names a nonnegative real \(a_v\), with tail at most \(2^{-s}\).
It is zero exactly when the machine does not halt. The purported Puiseux
procedure would decide halting.
\end{proof}

This obstruction already concerns formal coefficients. Arbitrary choices of
ramification or convergence radius do not help. Every \(a_v\) is in fact a
dyadic rational; the missing information is its \emph{exact presentation}.
Merely knowing that Cauchy-named values happen to be algebraic is insufficient.
The proposition is a computational obstruction to the stated uniform output,
not a claim of unprovability for every possible statement called Puiseux's
theorem.

The proof below also works over another fixed effectively presented
characteristic-zero coefficient field if the exact decisions, splitting
operations, and quantitative bounds used below are provided with their
\ipra\ correctness proofs. This certificate-relative extension does not
assert that arbitrary complex names automatically supply those data.

\section{Arithmetic conventions and logical discipline}
\label{sec:arithmetic}
We use intuitionistic logic with equality, all primitive recursive defining
equations, successor separation and injectivity, and quantifier-free
induction. Integers, rationals, arrays, polynomials, and matrices are coded
by natural numbers. Every finite search has a bound supplied before it runs.
A bounded conjunction is a primitive recursive Boolean fold, whose assertion
is an equation such as \(b=1\). Fixing a finite horizon and proving its
table invariant therefore uses only quantifier-free induction; the horizon
is universally generalized afterwards.

Congruence modulo \(T^L\) means equality of coefficients
\(0,\ldots,L-1\). It is a finite decidable predicate.
Equality of series abbreviates the family of these predicates for all \(L\),
and is never used as a decidable test.
A \emph{witnessed valuation} is a verified zero prefix and a displayed
nonzero coefficient. We never search indefinitely for such a coefficient.

The discriminant and resultant are finite determinant polynomials. For
monic polynomials we use
\begin{align}
 \disc(G(Y+c))&=\disc(G(Y)),\label{eq:disc-translation}\\
 \disc(T^{-pr}G(T^q,T^pY))
   &=T^{-pr(r-1)}\disc(G)(T^q),\label{eq:disc-scale}\\
 \disc\Bigl(\prod_jG_j\Bigr)
   &=\prod_j\disc(G_j)\prod_{i<j}\Res(G_i,G_j)^2.
   \label{eq:disc-product}
\end{align}
These are universal finite polynomial identities. One proof expands them
with finitely many indeterminate roots and then eliminates symmetric
monomials by the elementary symmetric-polynomial algorithm; degrees bound
the finite monomial list. Equivalently use the determinant identities.
Neither proof assumes a splitting theorem for power series.

\section{Squarefree reduction and one polynomial denominator}
\label{sec:reduction}
Translate the exact centre to \(X=0\), and put \(K=\A\) with the fixed
effective presentation proved in Section~\ref{sec:backend}.
Arithmetic in \(K(X)\) is finite rational-function arithmetic.
Compute the characteristic-zero squarefree decomposition
\begin{equation}
 \frac{P(z_0+X,Y)}{a_d(X)}=\prod_{m=1}^d Q_m(X,Y)^m.
 \label{eq:squarefree}
\end{equation}
Nonconstant \(Q_m\) are monic, squarefree, and pairwise coprime in \(K(X)[Y]\).
Empty factors are \(1\).

Here is a gcd-only algorithm and its verification. For a nonconstant monic
\(Q\), put \(G=\gcd(Q,Q_Y)\), \(S=Q/G\). Then \(S\) is squarefree.
To prove this without irreducible factorization, set
\(A=\gcd(S,S_Y)\), \(D=\gcd(A,A_Y)\), and write
\(A=uD\), \(A_Y=vD\), \(S=wA\). Since \(A\mid S_Y\),
coprimality of \(u,v\) implies \(u\mid w\).
Also \(A\mid Q,Q_Y\), hence \(A\mid G\).
If \(A^k\mid G\), differentiation gives
\(DA^{k-1}\mid G_Y\), and therefore \(A^{k+1}\mid SG_Y\).
Since \(A\mid S_Y\), we also have \(A^{k+1}\mid S_YG\).
Thus \(A^{k+1}\mid Q_Y\) and \(A^{k+1}\mid Q=SG\), so
\(A^{k+1}\mid G\). Induction only through \(\deg G+1\) forces \(\deg A=0\).
This proves the squarefreeness assertion by bounded finite algebra.

Repeatedly output \(S=Q/\gcd(Q,Q_Y)\) and replace \(Q\) by its gcd.
In characteristic zero the gcd degree is strictly smaller than that of a
nonconstant \(Q\), so after at most \(d\) stages this expresses the original
monic polynomial as a product of squarefree factors. To group their
overlaps, maintain pairwise coprime squarefree factors with multiplicities.
When inserting a new squarefree factor \(S\), scan the pre-insertion list.
For each factor \(A\)
of multiplicity \(m\), set \(D=\gcd(A,S)\), replace \(A\) by \(A/D\)
with multiplicity \(m\) and \(D\) with multiplicity \(m+1\), and replace
\(S\) by \(S/D\). After that scan, append the remaining \(S\), if nonconstant,
with multiplicity \(1\). Ignore factors equal to \(1\). Squarefreeness implies
\(\gcd(D,A/D)=\gcd(D,S/D)=1\); Bezout identities preserve coprimality with
the other factors. The product with multiplicities is preserved at every
step. The list has at most \(d\) nonconstant entries, so at most \(d^2\)
such refinements suffice. Group entries of the same multiplicity to obtain
the \(Q_m\). All squarefreeness and coprimality claims follow by the displayed
gcd argument and finite Bezout calculations.

Choose a nonzero polynomial \(C(X)\) divisible by every coefficient
denominator in every \(Q_m\), for instance their product.
If \(n_m=\deg_YQ_m\), put
\[
 F_m(X,Z)=C(X)^{n_m}Q_m(X,Z/C(X)),\qquad F(X,Z)=\prod_mF_m(X,Z).
\]
The coefficient at \(Z^{n_m-i}\) is \(C^i\) times the corresponding
coefficient of \(Q_m\), hence is polynomial. Thus the \(F_m\) are monic
polynomials in \(K[X,Z]\), squarefree and pairwise coprime over \(K(X)\).
Their product \(F\) is monic and squarefree. Define
\[
 n=\deg_ZF=\sum_m n_m\le d,\qquad \delta=\ord_X\disc_Z(F).
\]
The discriminant is a nonzero finite polynomial; its valuation is a finite
scan. A monic linear polynomial has discriminant \(1\).
By \eqref{eq:disc-product}, \(\ord\disc(F_m)\le\delta\), since all
discriminants and resultants here have nonnegative valuations.

Write
\begin{equation}
 C(X)=X^vU(X),\qquad U(0)\ne0.
 \label{eq:C-unit}
\end{equation}
A root \(Z_i(T)\) of \(F_m(T^N,Z)\) yields
\[
 Y_i(T)=T^{-Nv}Z_i(T)/U(T^N).
\]
Thus \(p=Nv\) is a common pole bound. The reciprocal of the polynomial
unit will be constructed explicitly.

\section{A bounded Newton--Hensel construction}
\label{sec:newton}
The degree-decreasing splitting idea has a constructive precedent in
Mannaa--Coquand \cite{mannaa}. We give the necessary bounded coefficient
searches and finite precision bookkeeping explicitly: constructivity alone
would not prove a primitive recursive bound.
Set
\begin{equation}
 N=(d!)^d,\qquad \Delta=N\delta,\qquad B=dN(\delta+1).
 \label{eq:budgets}
\end{equation}
For explanatory notation we describe coefficient streams, but the algorithm
will operate only on finite arrays.

\subsection{Coprime Hensel lifting is a finite recurrence}
\begin{lemma}[Coprime lifting]
\label{lem:hensel}
Suppose \(G(T,Z)\) is monic of degree \(r\) and
\(G(0,Z)=\prod_{j=1}^s C_j(Z)\), with pairwise coprime monic \(C_j\)
of degrees \(r_j\).
There are unique coefficients \(V_{j,\ell}\), \(\deg_ZV_{j,\ell}<r_j\), in
\[
 G_j(T,Z)=C_j(Z)+\sum_{\ell\ge1}T^\ell V_{j,\ell}(Z)
\]
whose product is \(G\) at every finite precision.
They are given by a primitive recursive coefficient recurrence.
\end{lemma}
\begin{proof}
Set \(D_j=\prod_{k\ne j}C_k\). Extended Euclid returns \(E_j\) with
\(E_jD_j\equiv1\pmod{C_j}\).
After computing all orders below \(\ell\), let \(R_\ell\) be the
coefficient of \(T^\ell\) in \(G-\prod_jG_j^{<\ell}\). It has degree
less than \(r\). Define
\[
 V_{j,\ell}=\operatorname{rem}(E_jR_\ell,C_j).
\]
The polynomial \(\sum_j V_{j,\ell}D_j-R_\ell\) has zero remainder modulo
each \(C_j\). Bezout identities imply divisibility by their product.
Its degree is less than \(r\), so it is zero.
This proves correctness at order \(\ell\); the same remainder calculation
proves uniqueness. No coefficient beyond order \(\ell\) is used.
\end{proof}

\subsection{Bounded slope selection and strict degree decrease}
At a node, \(G(X,Y)\) is monic integral of degree \(r\ge2\), with
witnessed discriminant valuation at most \(e\delta\), where \(e\) is
the product of earlier ramifications. Translate
\(Y=Z-a_1(X)/r\) to remove the next-to-leading coefficient:
\[
 \widetilde G(X,Z)=Z^r+\sum_{i=2}^r b_i(X)Z^{r-i}.
\]
Some \(b_i\) has a nonzero coefficient at index at most \(e\delta\).
Otherwise \(\widetilde G\equiv Z^r\pmod{X^{e\delta+1}}\); the discriminant
polynomial would vanish through that index, a contradiction.

Inspect every \(b_i\) through index \(B\), and choose the least ratio \(j/i\)
among discovered first nonzero coefficients. This is the true least slope:
a known ratio is at most \(e\delta\le\Delta\), whereas any undiscovered
coefficient has every index at most \(B\) zero and any later ratio exceeds
\(B/r\ge N(\delta+1)>\Delta\).
No test of infinite-series zero is being made.
Write the ratio as \(p_0/q\) in lowest terms. Then
\[
 0\le p_0\le B,\quad 1\le q\le r,\quad q\mid d!.
\]
The transform
\begin{equation}
 G^*(T,W)=T^{-p_0r}\widetilde G(T^q,T^{p_0}W)
 \label{eq:newton-transform}
\end{equation}
is monic integral. Its coefficient of \(W^{r-1}\) is zero; its
specialization \(D(W)=G^*(0,W)\) has a nonzero lower coefficient, because
a slope minimizer attains equality.

It is not \((W-c)^r\): the next coefficient would give \(-rc=0\), hence
\(c=0\), contrary to \(D\ne W^r\).
The exact splitting algorithm of Section~\ref{sec:backend} gives
\[
 D(W)=\prod_c(W-c)^{r_c},
 \quad \#\{c\}\ge2,\quad 1\le r_c<r,\quad \sum_c r_c=r.
\]
Apply Lemma~\ref{lem:hensel} to these coprime groups and recurse on their
factors. At degree one the root is the negative constant coefficient.

\subsection{Discriminant invariant, ramification, and termination}
Degree strictly decreases along each path. The forest has at most \(d-1\)
internal nodes and at most \(2d-1\) nodes altogether.
Each local \(q\) divides \(d!\); there are at most \(d-1\) such steps per
path, so each accumulated ramification divides \(N=(d!)^d\).
This divisibility justifies replacing every leaf parameter by a power of
one common \(T\) with \(X=T^N\).

Translation leaves the discriminant unchanged, while
\[
 \ord\disc(G^*)=q\ord\disc(G)-p_0r(r-1)\le qe\delta.
\]
The Hensel factors have pairwise resultants with nonzero constant terms.
Their discriminant valuations therefore add to that of \(G^*\), by
\eqref{eq:disc-product}. Each child has a witnessed valuation at most
\(qe\delta\), which is exactly the invariant with the new path
ramification. The existence of these witnesses follows by a finite product
scan through the known bound. We do not decide whether a child
discriminant stream is zero.

Reconstructing a parent root uses the inverse affine transformation.
If the child parameter has \(T=S^f\), then
\[
 X=S^{qf},\qquad
 Y=-a_1(S^{qf})/r+S^{p_0f}Z_{\rm child}(S).
\]
All exponents remain nonnegative. The child degrees sum to the parent
degree, so the leaf count for \(F_m\) is \(n_m\).

\begin{remark}
The common denominator is deliberately generous. The tempting bound
\(q r_c\le r\) fails for a zero characteristic root:
\(W^3(W^2+1)\) has \(r=5,q=2,r_0=3\).
The proof does not use that claim. Smaller certified common ramifications
can be used in examples without improving the universal bound.
\end{remark}

\subsection{An explicit global precision budget}
\label{subsec:precision}
For output precision \(L\), begin with input arrays through
\begin{equation}
 J(L)=\max(L,B)+d^2B+2.
 \label{eq:lookahead}
\end{equation}
Translation and Hensel lifting consume no extra coefficient positions.
In \eqref{eq:newton-transform}, an output coefficient of index \(k\)
requires an input coefficient of index at most \(k+p_0r\le k+dB\).
Fewer than \(d\) levels consume at most \(d^2B\); the maximum with \(B\)
covers each slope scan.

Run the forest in at most \(d\) rounds, splitting all nonterminal nodes,
with the remaining array length reduced by \(dB\) per level.
Node counts, coefficient loops, and all polynomial and exact-field
operations are explicitly bounded.
Reconstruction and \(T\mapsto T^{N/e}\) require no larger index.
This gives a finite-array primitive recursive program, rather than an
unsupported assertion about nested recursion on arbitrary series evaluators.
Deterministic splitting choices and Hensel uniqueness give compatible
prefixes when the requested precision is increased.

\section{One finite jet gives polynomial simple-root equations}
\label{sec:bootstrap}
Define
\begin{equation}
 h=\Delta+1,\qquad M=2h+1.
 \label{eq:bootstrap-budget}
\end{equation}
Run the finite construction once, at precision \(M\). Its output polynomials
\(b_i(T)\), truncated below \(M\), satisfy
\begin{equation}
 F_m(T^N,Z)\equiv\prod_{i\in I_m}(Z-b_i(T))\pmod{T^M}.
 \label{eq:jet-splitting}
\end{equation}
The group \(I_m\) retains the multiplicity label \(m\).

\subsection{Finite discriminant congruence gives every needed witness}
Multiply the group identities and apply the discriminant polynomial:
\begin{equation}
 \disc(F)(T^N)\equiv\prod_{i<j}(b_i-b_j)^2\pmod{T^M}.
 \label{eq:finite-disc}
\end{equation}
Its left side has witnessed leading index \(\Delta<M\).
Each difference consequently has a first nonzero coefficient at an index
\(k_{ij}\le\Delta\). Otherwise the product would vanish through \(\Delta\).
Multiplication of the leading terms gives
\[
 2\sum_{i<j}k_{ij}=\Delta.
\]
For \(n=1\) the product is empty and \(\Delta=0\).
Differentiating the group identity and evaluating at \(b_i\) gives
a witnessed derivative order
\[
 s_i=\ord F_{m,Z}(T^N,b_i)
     =\sum_{\substack{j\in I_m\\j\ne i}}k_{ij}\le\Delta<h.
\]
All of these numbers are obtained from finite arrays.

\subsection{Desingularization by a finite polynomial shift}
Let \(a_i(T)\) be the truncation of \(b_i\) \emph{through degree \(h\)};
then \(T^{h+1}\mid(b_i-a_i)\). Put
\begin{equation}
 H_i(T,V)=T^{-(h+s_i)}
       F_m(T^N,a_i(T)+T^hV).
 \label{eq:H}
\end{equation}
\begin{lemma}[Finite desingularization]
\label{lem:desingularization}
This expression is a polynomial, with
\(H_i(0,0)=0\) and \(H_{i,V}(0,0)\ne0\).
\end{lemma}
\begin{proof}
The root jet gives \(T^M\mid F_m(T^N,b_i)\).
Finite Taylor expansion shows that the derivative at \(a_i\) differs from
the derivative at \(b_i\) by a multiple of \(T^{h+1}\), preserving its
leading order \(s_i\). The same expansion gives
\[
 T^{h+s_i+1}\mid F_m(T^N,a_i):
\]
the linear correction has that divisibility, the higher corrections have
order at least \(2h+2\), and \(M=2h+1\ge h+s_i+1\).

After division in \eqref{eq:H}, the constant-in-\(V\) term is divisible
by \(T\). The linear coefficient is
\(T^{-s_i}F_{m,Z}(T^N,a_i)\), with a nonzero constant term.
For degree \(j\ge2\) in \(V\), the remaining factor is
\(T^{(j-1)h-s_i}\), a positive power. Hence all exponents are nonnegative
and both conclusions hold.
\end{proof}
The truncation convention matters: cutting only below \(h\) could leave a
nonzero value \(H_i(0,0)\). Including the coefficient of degree \(h\)
avoids a further constant translation.
The Newton forest is now finished. All later coefficients will be produced
from these finitely many polynomial simple-root equations.

\section{A coefficient recurrence and an explicit convergence proof}
\label{sec:analytic}

The finite algebraic preparation has supplied polynomial equations with a
simple root at the origin. We now construct their entire coefficient sequences
and prove convergence by estimates on finite polynomials. In particular, no
analytic implicit-function theorem is being assumed.

\begin{lemma}[A quantitative simple-root construction]\label{lem:analytic}
Let $H(T,V)\in\A[T,V]$ satisfy
$H(0,0)=0$ and $\lambda=H_V(0,0)\ne0$. Put
\[
 \Phi(T,V)=V-\lambda^{-1}H(T,V)=\sum_{i,j}b_{ij}T^iV^j.
\]
Choose rational bounds $\beta_{ij}\ge\abs{b_{ij}}$, with
$\beta_{00}=\beta_{01}=0$, and a rational $S$ such that
\begin{equation}
 S\ge1,\qquad S\ge\sum_{i,j}(j+1)\beta_{ij},\qquad
 r=\frac1{4S},\qquad \rho=r^2=\frac1{16S^2}.
 \label{eq:radius}
\end{equation}
There is a unique zero-constant formal solution
$u(T)=\sum_{n\ge1}u_nT^n$ of $H(T,u(T))=0$. Its coefficients are given by
\begin{equation}
 u_n=[T^n]\Phi\left(T,\sum_{k=1}^{n-1}u_kT^k\right)
 \qquad(n\ge1).
 \label{eq:implicit-recursion}
\end{equation}
For every finite $L$,
\begin{equation}
 \sum_{n=1}^{L}\abs{u_n}\rho^n\le r.
 \label{eq:weighted}
\end{equation}
Consequently the series converges for $\abs t<\rho$. If
$u_L(t)=\sum_{n=1}^{L}u_nt^n$ and $\abs t\le q\rho$, $0\le q<1$, then
\begin{align}
 \abs{u_M(t)-u_L(t)}&\le r q^{L+1} &&(M>L),\\
 \abs{\lambda^{-1}H(t,u_L(t))}
       &\le\frac{5r}{4}q^{L+1}.\label{eq:residual}
\end{align}
For $\abs t\le\rho$, there is at most one root $v$ of $H(t,v)=0$
with $\abs v\le r$.
\end{lemma}

\begin{proof}
Normalization gives $b_{00}=b_{01}=0$. Thus every monomial of $\Phi$
has either $i\ge1$ or $j\ge2$. When a zero-constant polynomial is
substituted for $V$, its coefficient of degree $n$ uses only coefficients
of that polynomial of degrees less than $n$. This proves that
\eqref{eq:implicit-recursion} is well defined and uniquely forces every
coefficient of the required identity. It is a primitive recursion with
finite coefficient convolutions.

Here is a finite proof of the bound. Define the nonnegative rational
polynomial $\Phi_+(T,W)=\sum\beta_{ij}T^iW^j$ and the Picard polynomials
\[
 X_0=W_0=0,\qquad X_{j+1}=\Phi(T,X_j),\qquad
 W_{j+1}=\Phi_+(T,W_j).
\]
Finite convolution inequalities show coefficientwise that
$\abs{[T^n]X_j}\le[T^n]W_j$. For $0\le w\le r$, every allowed
monomial $\rho^iw^j$ is at most $\rho$: use $\rho,w\le1$ when
$i\ge1$, and $w^j\le r^2=\rho$ when $i=0$, $j\ge2$. Therefore
\[
 \Phi_+(\rho,w)\le\rho\sum\beta_{ij}\le S\rho=r/4.
\]
Induction gives $W_j(\rho)\le r$ for every finite $j$.

If two zero-constant polynomials agree modulo $T^k$, $k\ge1$,
their images under $\Phi$ agree modulo $T^{k+1}$. Indeed, factor
$V^j-W^j=(V-W)\sum_{\ell<j}V^\ell W^{j-1-\ell}$: either the
factor $T^i$ or one of the additional zero-constant factors supplies the
extra power of $T$. Consequently $X_{j+1}-X_j$ is divisible by
$T^{j+1}$, and the first $L$ coefficients of $X_L$ are exactly the
coefficients in \eqref{eq:implicit-recursion}. Their weighted absolute
sum is at most $W_L(\rho)\le r$, proving \eqref{eq:weighted}.
All Picard calculations may be truncated modulo $T^{L+1}$.

For $M>L$ and $\abs t\le q\rho$, the finite tail is bounded by
\[
 \sum_{n=L+1}^{M}\abs{u_n}\rho^n q^n
 \le q^{L+1}\sum_{n=L+1}^{M}\abs{u_n}\rho^n
 \le r q^{L+1}.
\]
This proves the Cauchy estimate without first assuming the existence of
an infinite sum. In particular, on $\abs t\le\rho/2$, truncation at
$L$ has error at most $r2^{-L-1}$.

For a finite polynomial $A$, write
$\norm A_\rho=\sum_n\abs{[T^n]A}\rho^n$.
Finite convolution gives $\norm{AB}_\rho\le\norm A_\rho\norm B_\rho$.
The polynomial
$R_L=u_L-\Phi(T,u_L)=\lambda^{-1}H(T,u_L)$ is divisible by
$T^{L+1}$ and satisfies
\[
 \norm{R_L}_\rho
 \le\norm{u_L}_\rho+\Phi_+(\rho,\norm{u_L}_\rho)
 \le r+r/4.
\]
Factoring $q^{L+1}$ from its finite coefficient sum proves
\eqref{eq:residual}. Thus the limit solves the equation with an explicit
residual modulus.

Finally, for $\abs u,\abs v\le r$ and $\abs t\le\rho$,
\[
 \abs{\Phi(t,u)-\Phi(t,v)}
 \le\sum_{\substack{i\ge0\\j\ge1}}\beta_{ij}\rho^i j r^{j-1}\abs{u-v}
 \le Sr\abs{u-v}=\tfrac14\abs{u-v}.
\]
In this sum, $\rho^ir^{j-1}\le r$: either $i\ge1$, or $i=0$
and $j\ge2$. Two fixed points are therefore equal. More quantitatively,
\[
 \abs{u-v}\le\tfrac43\left(
 \abs{\lambda^{-1}H(t,u)}+\abs{\lambda^{-1}H(t,v)}\right).
\]
This also gives uniqueness directly in terms of rational approximation
bounds, without a decision procedure for equality of represented complex
numbers.
\end{proof}

The infinite-series notation in this lemma has a precise arithmetic
meaning: the coefficient program is \eqref{eq:implicit-recursion}, and
its asserted properties are the displayed finite inequalities for arbitrary
finite indices. Each coefficient has a finite exact algebraic code.
Evaluation to a requested precision uses certified rational approximations
and a finite error budget for the finitely many terms. No completeness or
choice axiom is added to $\ipra$.

\subsection{Recovering the full factorization with a finite precision budget}

Apply the lemma to every prepared polynomial $H_i$. Write
\[
 Z_i(T)=a_i(T)+T^h u_i(T)
\]
for the resulting formal root of the corresponding
$F_m(T^N,Z)$. The finite construction gave jets $b_i$ modulo $T^M$,
where $M=2h+1$ and $h=\Delta+1$, together with witnessed first
differences
\[
 k_{ij}=\ord_T(b_i-b_j),\qquad
 2\sum_{i<j}k_{ij}=\Delta.
\]
The new root agrees with $b_i$ through degree $h$, because $u_i(0)=0$.
It therefore preserves the known nonzero coefficient at degree $k_{ij}$.
This proves distinctness by finite witnesses, rather than by searching
indefinitely for a nonzero difference.

For completeness, the following argument supplies an explicit precision
budget for recovering the splitting identity. Fix one squarefree group,
let $I_m$ index its $n_m$ roots, and set
\[
 R(T,Z)=F_m(T^N,Z)-\prod_{i\in I_m}(Z-Z_i(T)).
\]
Both polynomials are monic of degree $n_m$, so $\deg_ZR<n_m$.
The Vandermonde matrix with rows
$(1,Z_i,\ldots,Z_i^{n_m-1})$ has determinant
$\prod_{i<j,\ i,j\in I_m}(Z_j-Z_i)$, with a witnessed nonzero
leading coefficient at degree
$v_m=\sum_{i<j,\ i,j\in I_m}k_{ij}\le\Delta$.
The adjugate identity expresses this determinant times each coefficient
of $R$ as a finite linear combination of the values $R(T,Z_i(T))$.

We use only a finite cancellation lemma. If $A$ has a nonzero leading
coefficient $a_v$ at degree $v$ and $AB\equiv0\pmod{T^L}$, then
$B\equiv0\pmod{T^{L-v}}$. At degree $v+k$, previous vanishing
coefficients of $B$ leave just $a_vb_k=0$; multiply by $a_v^{-1}$
and induct on $k$. Thus, to prove factorization through degree $K$, it
suffices to compute the root identities modulo $T^L$ for
$L=K+\Delta+1$. The adjugate identity and finite cancellation then
make every coefficient of $R$ vanish through degree $K$. Since $K$
was arbitrary, we obtain
\begin{equation}
 F_m(T^N,Z)=\prod_{i\in I_m}(Z-Z_i(T)).
 \label{eq:exact-factor}
\end{equation}
This proves the formal identity and its finite certificate requirements.
The convergence estimates just proved also justify its evaluated version:
the finite identity
\[
 \prod_j a_j-\prod_j b_j
 =\sum_j(a_j-b_j)\prod_{\ell<j}a_\ell\prod_{\ell>j}b_\ell
\]
bounds the difference of two products of $n$ factors of magnitude at most
$B$ by $nB^{n-1}\max_j\abs{a_j-b_j}$. Apply this with the explicit
branch bounds and tail moduli. Thus every evaluated truncation error has
a stated finite bound tending to zero.

\section{Units, poles, and all branches}
\label{sec:laurent}

\subsection{An explicit reciprocal and one common radius}

Recall the common denominator $C(X)=X^vU(X)$, where $U(0)\ne0$,
and put $p=Nv$. Division by $C(T^N)$ introduces only the finite pole
$T^{-p}$ and the reciprocal of a polynomial with nonzero constant term.
Here is a direct coefficient construction of that reciprocal.

Write $B(T)=U(T^N)=b_0(1+V(T))$, where $V(0)=0$.
Choose rational numbers $A\ge\max(1,\abs{b_0^{-1}})$ and
$L\ge1+\sum_{j\ge1}\abs{[T^j]V}$, using rational upper bounds
for the finitely many algebraic coefficients, and set $R_u=1/(2L)$.
Then $\norm V_{R_u}\le1/2$. Define
\[
 d_0=b_0^{-1},\qquad
 d_n=-b_0^{-1}\sum_{j=1}^{\min(n,\deg B)}b_jd_{n-j}.
\]
The first $K+1$ coefficients of this reciprocal agree with those of
$b_0^{-1}\sum_{j=0}^{K}(-V)^j$, because $V$ has zero constant term.
Hence, for every finite $K$,
\[
 \sum_{n=0}^{K}\abs{d_n}R_u^n
 \le A\sum_{j=0}^{K}(1/2)^j\le2A.
\]
Coefficient comparison proves
$B(T)\sum_{n\ge0}d_nT^n=1$, and the same finite tail argument as
before proves convergence on $\abs T<R_u$. On the smaller disc
$\abs T\le R_u/2$, the tail beyond $K$ is at most $2A2^{-K-1}$.
Also $\abs{V(T)}\le1/2$ for $\abs T\le R_u$, so
$\abs{B(T)}\ge\abs{b_0}/2$: the reciprocal is an actual local
inverse, with an explicit apartness bound.

Choose a positive rational $\eta\le1$ no larger than $R_u$ and all
the implicit-equation radii $\rho_i$. If needed, shrink it by the same
unit estimate for the original leading coefficient $a_d(T^N)$ after
removing its initial power of $T$. Thus this coefficient is apart from
zero at every punctured argument under consideration.
Choose rational $\alpha_{ik}\ge\abs{[T^k]a_i}$ and put
\[
 B_i=1+\sum_k\alpha_{ik}\eta^k+\eta^h r_i,
 \qquad M_i=2AB_i.
\]
The preceding estimates give finite weighted bounds $B_i$ for $Z_i$
and $2A$ for $B^{-1}$. Finite convolution therefore gives, for
\[
 W_i(T)=\frac{Z_i(T)}{U(T^N)}=\sum_{n\ge0}w_{i,n}T^n,
 \qquad Y_i(T)=T^{-p}W_i(T),
\]
the estimate $\sum_{n=0}^{K}\abs{w_{i,n}}\eta^n\le M_i$ for every
finite $K$. All $w_{i,n}$ are obtained by the displayed primitive
recursions and finite convolutions.

On a rationally specified annulus
$0<\sigma\le\abs T\le\eta/2$, we consequently have
\begin{equation}
 \left|Y_i(T)-T^{-p}\sum_{n=0}^{K}w_{i,n}T^n\right|
 \le \sigma^{-p}M_i\,2^{-K-1}.
 \label{eq:annulus-tail}
\end{equation}
For example, $K=k+\lceil\sigma^{-p}M_i\rceil+1$ more than suffices
for truncation error $2^{-k}$. This is a primitive recursive modulus;
finite coefficient-evaluation errors receive an additional explicit portion
of the error budget. The positive lower bound $\sigma$ is the constructive
information needed to evaluate negative powers.
With dyadic input bounds \(2^{-b}\le\abs T\) and \(M_i\le2^{\mu_i}\),
the sharper choice \(K=k+bp+\mu_i+2\) gives truncation error at most
\(2^{-k-3}\). The integer \(\mu_i\) is computed from a rational numerator
bit bound; no logarithm or unbounded search is needed.

If $w_{i,0},\ldots,w_{i,p-1}$ all vanish, the branch has a finite
limit $w_{i,p}$ and an ordinary power-series parametrization after
$X=T^N$. Otherwise the first nonzero coefficient among these finitely
many gives its pole. Exact algebraic equality makes this a bounded test.
No decision about whether an entire infinite series vanishes is required.

\subsection{Exhaustiveness, multiplicities, and constructive root selection}

Undoing the monic denominator substitution in
\eqref{eq:exact-factor} and restoring multiplicities gives
\[
 P(z_0+T^N,Y)=a_d(T^N)\prod_i\bigl(Y-Y_i(T)\bigr)^{m_i}.
\]
This is the complete factorization, including multiplicities, on the
punctured disc. Together with \eqref{eq:annulus-tail} it supplies the
required convergent Puiseux expansions at the original base point \(z_0\).

A pointwise statement selecting the index of a represented root needs one
more argument. Constructively, a zero product of arbitrary represented
complex numbers does not by itself justify a disjunction of zero factors.
Here the finite preparation supplies the missing separation information.
For two distinct branches, multiplication by the common reciprocal unit
preserves the known first difference $k_{ij}$. Choose a positive rational
$\ell_{ij}$ below the magnitude of that coefficient in $W_i-W_j$.
The weighted bound gives, for $\abs T\le\eta$,
\[
 \left|\sum_{n>k_{ij}}(w_{i,n}-w_{j,n})T^{n-k_{ij}}\right|
 \le(M_i+M_j)\abs T\,\eta^{-k_{ij}-1}.
\]
Shrink the radius to the positive rational
\[
 \eta_* =\min\left(\frac\eta2,
  \min_{i<j}\frac{\ell_{ij}\eta^{k_{ij}+1}}{2(M_i+M_j)}\right),
\]
omitting the pairwise minimum when there is only one distinct branch.
Then
\[
 \abs{Y_i(T)-Y_j(T)}
 \ge\frac{\ell_{ij}}2\abs T^{k_{ij}-p}
 \qquad(0<\abs T\le\eta_*).
\]
On $\sigma\le\abs T\le\eta_*\le1$, a positive rational lower
bound for this distance is
$(\ell_{ij}/2)\sigma^{\max(k_{ij}-p,0)}$.
Take the minimum of these finitely many lower bounds and call it $\delta$.

Given a represented root $y$, approximate each $y-Y_i(T)$ with error
at most $\delta/16$, and test whether the squared norm of its rational
approximation is less than $\delta^2/9$. Some test must succeed:
otherwise every exact factor has magnitude at least $13\delta/48$,
contradicting the factorization and the apart leading coefficient.
There is at most one successful test, since two corresponding exact
distances would have sum less than $38\delta/48<\delta$.
For the selected index all other factors are bounded away from zero
by $29\delta/48$. Cancel them; the remaining equation
$(y-Y_i(T))^{m_i}=0$ implies $y=Y_i(T)$ by the norm identity.
Thus root selection is a bounded rational approximation procedure.
With a single distinct branch, the final power equation gives the same
conclusion directly. This establishes pointwise exhaustiveness without
excluded middle for complex equality.


\section{The precise first-order \ipra\ conclusion}
\label{sec:internal}
The internal assertion does not quantify over an infinite set of series.
Let \(\operatorname{Valid}(p)\) check the exact finite input.
The fixed program \(\operatorname{Build}(p)\) returns the finite packet
\[
 (N,p_{\rm pole},\{m_i,a_i,h,s_i,H_i\}_i,
 U,u_0,\eta,\{M_i\}_i,\text{finite certificates}).
\]
The fixed program \(\operatorname{Jets}(p,L)\) uses the simple-root
recurrence, the unit reciprocal recurrence, and finite convolution to return
the numerator coefficients through \(L\). Malformed inputs have specified
default outputs; validity proofs show that the defaults are not used.

\begin{center}
\small
\begin{tabular}{p{.24\linewidth}p{.34\linewidth}p{.31\linewidth}}
\toprule
Operation & Finite budget & Invariant\\
\midrule
Squarefree reduction & \(d\) stripping stages, at most \(d^2\) refinements & Quotients and Bezout identities\\
Newton forest & At most \(d\) rounds, \(2d-1\) nodes & Degree sums, path ramification\\
Slope scan & Indices through \(B\) & Discriminant witness at most \(e\delta\)\\
Initial jets & \(M=2(N\delta+1)+1\), input \(J(M)\) & Complete splitting congruences\\
Simple coefficients & Requested degree & Zero residual coefficients\\
Factorization & Residual horizon \(K+\Delta+1\) & Vandermonde and leading cancellation\\
Convergence & Picard iterates through a finite horizon & Weighted coefficient bounds\\
Numerical values & Explicit tail and evaluation precisions & Cauchy and residual estimates\\
\bottomrule
\end{tabular}
\end{center}

Finite sums of exact algebraic moduli have exact algebraic order tests.
Alternatively all majorants may use rational coordinate upper bounds.
An inequality at a represented complex argument denotes its rational
approximation version with a specified error budget.
A limiting equality denotes error bounds by \(2^{-k}\) for every
natural precision \(k\).

The actual formulas include
\[
 \forall p,L\;\bigl(\operatorname{Valid}(p)=1\ \longrightarrow\
 \operatorname{FactorCheck}(p,\operatorname{Build}(p),L,
                         \operatorname{Jets}(p,L))=1\bigr),
\]
and analogous finite prefix, weighted-bound, and tail assertions.
Their checkers are defined by the identities and inequalities proved in
this report, not by predicates whose mathematical meaning is assumed.
All existential outputs are witnessed by the displayed fixed procedures.

For each induction, fix a horizon and collect all already established
bounded coefficient identities and code bounds into a Boolean table
predicate. Its equation is a quantifier-free induction formula.
This applies to Euclidean division, determinant expansion, Hensel
corrections, stabilization, convolution bounds, and leading cancellation.
Afterwards universally generalize over the horizon.
This is the promised nonformal proof in \ipra; constructing its explicit
object derivation in Lean is a separate task described below.

No universal primitive recursive evaluator for arbitrary primitive
recursive programs is assumed. The backend, constructor, and jet evaluator
are fixed programs, with finite codes as numerical parameters.
The reused FTA theorem is instantiated with one fixed coefficient evaluator
for the exact algebraic code representation.

\section{Exact examples and executable evidence}
\label{sec:examples}
The Python program \texttt{code/check\_examples.py} uses Gaussian rationals
as pairs of exact \texttt{Fraction}s and only the standard library.
It implements the simple-root recurrence, finite bivariate polynomial
arithmetic, weighted bounds, and complete products with multiplicities.
It is a finite-jet checker and illustration, not a general Puiseux solver
or a formal proof of the theorem.

\begin{center}
\small
\begin{tabular}{p{.35\linewidth}p{.10\linewidth}p{.44\linewidth}}
\toprule
Polynomial & \(X=t^e\) & Branches or leading terms\\
\midrule
\(Y^2-X^3\) & \(e=2\) & \(\pm t^3\)\\
\((Y^2-X^3)^2-X^7Y\) & \(e=4\) &
\(s t^6+s c t^{11}-s c^3t^{21}/8+\cdots,\quad c^2=s/4\)\\
\((Y^2-X^3)^3(Y-X^2)\) & \(e=2\) &
\(\pm t^3\) each with multiplicity \(3\); \(t^4\) with multiplicity \(1\)\\
\(XY^2-(1+X)\) & \(e=2\) &
\(\pm t^{-1}(1+t^2/2-t^4/8+\cdots)\)\\
\((Y-X^{20})^2-X^{41}\) & \(e=2\) &
\(t^{40}\pm t^{41}\)\\
\bottomrule
\end{tabular}
\end{center}
In row two take \(s=1,c=\pm1/2\) and \(s=-1,c=\pm i/2\). The exact equation
for the correction is
\[
 (2V+t^5V^2)^2-s(1+t^5V)=0,\qquad V(0)=c.
\]
Its derivative is \(8c\ne0\). The exponents \(6\) and \(11\) give genuine
fourth-root ramification. The last example has degree two but discriminant
\(4X^{41}\); degree alone cannot bound the initial singularity scan.

The recorded run checks through degree \(96\). It checks all five complete
factorizations, implicit residuals and weighted bounds where an infinite
series is used, and multiplicities. Laurent products retain extra high
coefficients before truncation to protect against negative powers.
The output is \texttt{evidence/example\_checks.json}.

\section{A clear path to complete Lean formalization}\label{sec:lean}

\subsection{The object-theory target}

The inspected ProveIt snapshot is \texttt{\pin}, using Lean~4.32.0.
Its \href{https://github.com/VladimirReshetnikov/ProveIt/tree/d851229c33ef395db8dc04baf28c6a291073f61e/Logic/PrimitiveRecursiveArithmetic/ConciseIPRA/Lean/ConciseIPRA}{\texttt{ConciseIPRA} sources}
already supply scoped syntax (\texttt{Syntax}), intuitionistic derivability
(\texttt{Deduction}), primitive-recursive function symbols and
quantifier-free induction (\texttt{PR}), and transport to the concise
native calculus (\texttt{NativeBiInterpretation}).
The required endpoint is a derivation \emph{inside that encoded theory}:
\begin{lstlisting}
open LeanProofs.ConciseIPRA
Prf PRAAx ([] : List (Fm PTm 0))
  (PuiseuxIPRA.PuiseuxSentence concretePrograms)
\end{lstlisting}
Here \texttt{concretePrograms} must contain the actual \texttt{PR} syntax
trees for input validation, packet construction, coefficient generation,
and the finite predicates of Section~\ref{sec:internal}.  Their decoding
theorems must establish the intended meanings; naming a Boolean program
\texttt{factorJet} does not establish its correctness.

The supplied \texttt{lean/PuiseuxTarget.lean} writes the closed formula
explicitly using \texttt{Fm.eq}, \texttt{Fm.and}, \texttt{Fm.imp}, and
\texttt{Fm.all}.  Its matrix asserts packet validity, the complete
denominator-cleared factorization identity, prefix compatibility, and
finite coefficient bounds; a second universally quantified matrix asserts
the rational tail estimate.  It also gives the type of the required
derivation and conditional transport of an already supplied derivation
through \texttt{pra\_iff\_native}.  This last transport does not construct
the missing derivation.

\subsection{Dependency contracts and implementation order}

The following sequence specifies concrete completion targets.  Every
algorithm operates on natural-number codes, with default outputs on
malformed data; valid-input proofs exclude those defaults.

\begin{enumerate}
\setlength{\itemsep}{3pt}
\item \textbf{Coded arithmetic.}
Implement signed integers, rational numbers with positive denominators,
finite lists, polynomial arrays, bounded scans and Boolean folds.
Prove decoding equations for addition, multiplication, comparison,
coefficient access and finite convolution.

\item \textbf{Exact algebraic coefficients.}
Formalize Section~\ref{sec:backend}: real-root isolation, algebraic
arithmetic, equality and order, resultant elimination and finite splitting.
The reused constructive FTA root-name construction requires its own
deep arithmetic derivation; the existing draft is not a compiled
FTA theorem.

\item \textbf{Squarefree preprocessing.}
Prove the degree-decreasing Euclidean algorithm and its product and
B\'ezout invariants.  Verify the multiplicities, common denominator \(C\),
monic polynomial factors \(F_m\), and the denominator-cleared identity
relating them to the original polynomial.

\item \textbf{Truncated polynomial algebra.}
Prove congruence preservation for multiplication, substitution,
ramification, translation and discriminants.  Prove the finite
Hasse--Taylor identity and guarded division by a monomial.
These identities determine the exact precision losses.

\item \textbf{Coprime Hensel recursion.}
Implement the correction
\(U_{j,k}=\operatorname{rem}(E_jR_k,C_j)\).
Prove the product identity through the next coefficient, uniqueness
at the current precision, and compatibility when the requested
precision is increased.

\item \textbf{Bounded Newton tree.}
Implement Section~\ref{sec:newton} as a bounded stack or \(d\) rounds,
with at most \(2d-1\) nodes.  Maintain degree accounting,
\(e\mid N\), the discriminant bound \(e\delta\), and the precision reserve.
The scan bound \(B\) and lookahead \(J(M)\) are part of the program.

\item \textbf{Finite discriminant separation.}
From the factorization modulo \(T^M\), prove the discriminant congruence,
find every first difference by a scan through \(\Delta\), and establish
\(2\sum_{i<j}w_{ij}=\Delta\).
Only equality of finite algebraic coefficients is tested.

\item \textbf{Simple polynomial leaves.}
Formalize Section~\ref{sec:bootstrap} with \(a_i\) through degree \(h\)
and \(M=2h+1\).  Prove \(s_i\le\Delta\), the divisibility defining \(H_i\),
\(H_i(0,0)=0\), and \(H_{i,U}(0,0)\ne0\).
All four conclusions are finite polynomial assertions.

\item \textbf{Coherent coefficients and complete factorization.}
Implement the simple coefficient recurrence and prove its finite
correctness and uniqueness.  Synthetic division and cancellation by
the recorded leading differences recover every factor \(F_m\).
Use residual congruences modulo \(T^{L+\Delta+1}\) when establishing
factorization through degree \(L\).

\item \textbf{Rational majorants and evaluation.}
Formalize Section~\ref{sec:analytic} with truncated positive Picard
polynomials, rational coefficient majorants and finite weighted sums.
Add unit inversion and convolution to obtain the common radius and
Laurent tail bound, including the supplied positive lower bound
for a nonzero argument.

\item \textbf{Internal proof assembly.}
Construct the programs from \texttt{PR.comp} and \texttt{PR.prec}.
For each loop, use \texttt{PRA.qfInd} on its finite Boolean invariant
with the horizon held as a parameter.  Universally generalize afterward
and assemble the displayed \texttt{Prf PRAAx} endpoint.

\item \textbf{Semantic interpretation and audit.}
Interpret the coherent coefficient and tail statements as convergent
Puiseux expansions.  Transfer the object derivation to the native
calculus if desired, inspect its type and \texttt{\#print axioms},
and rebuild its complete dependency closure from the pinned sources.
\end{enumerate}

\subsection{Current status and the remaining work}

No Lean executable was available in this run.  The supplied Lean source
is therefore an \emph{uncompiled interface and outline}, including its
small structural proof terms.  It contains no omitted proof term or
postulated Puiseux theorem, and supplies no inhabitant of the final
derivation type.  The exact-algebraic backend, quantitative FTA dependency,
Newton--Hensel invariants, majorant bounds and their object-theory
derivations remain substantial formalization tasks.

The accompanying \texttt{lean/LEAN\_PLAN.md} expands the contracts above,
including the exact original-polynomial identity checked by
\texttt{factorJet}.  The reference computations provide finite evidence
for executable components.  A conventional Lean theorem about
\texttt{Complex} or \texttt{PowerSeries} may help organize the mathematics,
but it does not replace the encoded \(\ipra\) derivation required here.



\appendix
\section{Exact algebraic coefficients: a bounded arithmetic construction}
\label{sec:backend}

This section supplies the exact coefficient field used in the main proof.
Its operations, equality tests, bounds, and polynomial splitting are
primitive recursive on finite codes. The only root-existence theorem reused
is the single-root constructive FTA described in Section~\ref{sec:fta}.
It is used once in the completeness argument for each splitting operation,
not as an oracle in the executed algorithm.

\subsection{Real-algebraic codes and a bounded isolation procedure}

A real-algebraic code is either a rational number, or a tuple
\[
 (p,a,b,A,B,R)
\]
of rational polynomials and rationals, with $R$ a positive integer.
The finite validity conditions require $p$ to be monic and nonconstant,
\[
 Ap+Bp'=1,\qquad
 R>1+\sum_{i<\deg p}\abs{p_i},\qquad
 -R\le a<b\le R,\qquad p(a)p(b)<0.
\]
Write
\[
 M=1+\sum_i\abs{B_i}R^i,\qquad
 L=1+\sum_{i=2}^{\deg p}\binom i2\abs{p_i}R^{i-2},
 \qquad \sigma=\frac1{2ML},
\]
and additionally require $b-a<\sigma$.
Every validity test is a finite rational calculation; in particular
squarefreeness is certified by the displayed B\'ezout identity.

The represented value is specified by rational bisection on $[a,b]$.
At an exact midpoint zero, all subsequent approximations equal that
midpoint. Otherwise retain the half interval with opposite endpoint signs.
After $s$ steps its length is $(b-a)2^{-s}$. Thus the midpoints have an
explicit Cauchy modulus. If
\[
 D=1+\sum_{i=1}^{\deg p}i\abs{p_i}R^{i-1},
\]
the finite identity for differences of powers gives
$\abs{p(x)-p(y)}\le D\abs{x-y}$ on $[-R,R]$.
The opposite endpoint signs then give
$\abs{p(c_s)}\le D(b-a)2^{-s}$ for the midpoint $c_s$.
These finite inequalities establish a root name without appealing to a
completeness axiom or to an already existing real root.

\begin{lemma}[Effective separation and isolation]
\label{lem:backend-isolation}
All real roots of a nonzero rational polynomial can be returned as a finite
list of valid real-algebraic codes, by a primitive recursive algorithm.
The list is complete even for roots initially supplied only by Cauchy names.
\end{lemma}

\begin{proof}
First compute the squarefree part by rational Euclidean division, and
normalize it to a monic polynomial $p$. Constant polynomials have no roots.
The finite squarefree calculation also certifies that the squarefree part
divides the original polynomial and that a power of the squarefree part,
with exponent at most the original degree, is divisible by the original
polynomial. These identities establish equality of the root sets without
using a disjunction principle for a product of arbitrary real numbers.
Extended Euclid gives $Ap+Bp'=1$. Choose an integer $R$ with the extra margin
$R>2+\sum_{i<\deg p}\abs{p_i}$, and form $M,L,\sigma,D$ as above.
The elementary geometric-sum root bound places all roots strictly in
$(-R,R)$.

At a root $\alpha$, the B\'ezout identity gives
\[
 \abs{p'(\alpha)}\ge\frac1M.
\]
For $\alpha,t\in[-R,R]$, the exact polynomial identity
\[
 p(t)-p(\alpha)-(t-\alpha)p'(\alpha)
       =(t-\alpha)^2 E(\alpha,t)
\]
has $\abs{E(\alpha,t)}\le L$. Indeed, the remainder for the monomial $X^i$
is
\[
 (t-\alpha)^2
 \sum_{k=0}^{i-2}(i-1-k)t^k\alpha^{i-2-k},
\]
and its coefficient sum is $\binom i2$.
If $\alpha,\beta$ are roots with $\abs{\alpha-\beta}\le\sigma$, this identity
implies
\[
 \frac{\abs{\alpha-\beta}}M
 \le L\abs{\alpha-\beta}^2
 \le \frac{\abs{\alpha-\beta}}{2M}.
\]
Hence $\alpha=\beta$. This argument does not cancel a possibly zero real.

Choose an integer $K>32ML$ and put $h=1/K<\sigma/16$.
Evaluate $p$ at every grid point $j/K$ in $[-R,R]$.
Return each exact grid zero as a rational code. For each cell with nonzero,
oppositely signed endpoint values, return its bisection code.
Its width is less than $\sigma$, so the preceding estimate certifies
uniqueness. A cell adjacent to a grid zero contains no other root.

For completeness, let $\alpha$ be any named root. If it lies strictly
between nonroot endpoints $l,u$ of a cell, then
\[
 p(t)=(t-\alpha)
       \bigl(p'(\alpha)+(t-\alpha)E(\alpha,t)\bigr).
\]
The expression in parentheses differs from $p'(\alpha)$ by less than
$1/(4M)$ throughout the cell. Consequently the endpoint values have
opposite signs, and local uniqueness identifies $\alpha$ with the
returned bisection root.

The selection of this cell does not assume decidable equality of an
arbitrary named real with a grid point. Compute the positive rational
\[
 \mu=\min_{\substack{g\text{ a grid point}\\p(g)\ne0}}
                         \frac{\abs{p(g)}}D.
\]
The list in this minimum is nonempty, since its two extreme grid points
are not roots. Every root has distance at least $\mu$ from every nonroot
grid point. Obtain a rational approximation $c$ to $\alpha$ with error
less than $\min(\mu,\sigma,h)/8$.
If $\abs{c-g}\le\sigma/2$ for a grid zero $g$, then
$\abs{\alpha-g}<\sigma$, and the local uniqueness estimate proves
$\alpha=g$. Otherwise $\alpha$ is positively separated from every grid
zero as well. The rational value $c$ now determines its cell, whose
endpoints are nonroots. This gives a bounded selection in all cases.

The list has no repetitions: different cells have disjoint interiors,
and a nonroot grid endpoint is positively separated from every root.
Distinct returned roots are apart. Cancelling their positively bounded
distance in the same remainder estimate gives separation at least
$1/(ML)$; the weaker bound $\sigma$ will suffice below.
All searches range over the predetermined $2RK+1$ grid points.
\end{proof}

\subsection{Arithmetic, equality, and quantitative bounds}

Let $\alpha,\beta$ have monic rational annihilators $p,q$.
Finite resultant identities give annihilators for their sum and product:
\[
 \Res_T\bigl(p(T),q(Z-T)\bigr),\qquad
 \Res_T\bigl(p(T),T^{\deg q}q(Z/T)\bigr).
\]
Both are monic in $Z$, of degree $(\deg p)(\deg q)$, and therefore nonzero.
Negation uses $p(-Z)$. For a certified nonzero $\alpha$, remove all powers
of $T$ from $p$, obtaining $p_0$ with $p_0(0)\ne0$; an annihilator for
$1/\alpha$ is
$Z^{\deg p_0}p_0(1/Z)$.
The latter polynomial is nonzero. The annihilation identities follow by
polynomial elimination and do not presuppose a complex splitting algorithm.

Apply Lemma~\ref{lem:backend-isolation} to the chosen annihilator.
Using the input bisection names, approximate the intended arithmetic result
within $h/4$, where $h$ is the new isolating grid's mesh.
Each candidate's rational midpoint is within $h/2$ of its root, while
distinct candidates are separated by more than $16h$.
Exactly one midpoint therefore lies within $h$ of the computed
approximation. Finite rational comparisons select its code.
For products, input root bounds give the required input precision.
For reciprocals, use the explicit positive lower bound established next.

\begin{lemma}[Decidable algebraic zero and sign]
\label{lem:backend-zero}
Zero, equality, and order are primitive recursive for real-algebraic codes.
A nonzero value comes with a positive rational lower bound on its absolute
value.
\end{lemma}

\begin{proof}
Let a named value $\gamma$ have a nonzero rational annihilator
$r(Z)=Z^k q(Z)$ with $q(0)\ne0$, obtained by removing the initial zero
coefficients. Put
\[
 \eta=\min\left(1,
       \frac{\abs{q(0)}}{1+\sum_{i\ge1}\abs{q_i}}\right)>0
\]
and compute a rational approximation $c$ with error at most $\eta/8$.
If $\abs c<\eta/2$, then $\abs\gamma<5\eta/8$, and
\[
 \abs{q(\gamma)-q(0)}
 \le\abs\gamma\sum_{i\ge1}\abs{q_i}
 <\frac58\abs{q(0)}.
\]
Thus $q(\gamma)$ is positively apart from zero. Since
$\gamma^k q(\gamma)=0$, it follows that $\gamma^k=0$ and $\gamma=0$.
This uses the elementary finite power estimate, not a
nonzero-to-apartness principle. If instead $\abs c\ge\eta/2$, then
$\abs\gamma\ge3\eta/8$, and $\gamma$ has the sign of $c$.
Both alternatives are decided by one rational comparison.
Apply this procedure to the difference of two codes to obtain equality
and order. The argument also applies directly to the named arithmetic
result before its output isolating code has been selected.
\end{proof}

Upper bounds are supplied by the root bound $R$.
Complex-algebraic codes are pairs of real-algebraic codes, representing
$a+ib$; their equivalence is decidable equality of both components.
The usual two-coordinate formulas give addition and multiplication.
For a nonzero pair, at least one component has a positive rational
absolute-value lower bound, which also bounds the complex modulus from
below. An upper bound is supplied by the sum of the two component bounds.
Inversion uses
\[
 (a+ib)^{-1}=\frac{a-ib}{a^2+b^2},
\]
whose denominator is certified positive by the same tests.
This constructs the represented field $\A$ and all operations required
in the main argument.

\subsection{A rational annihilator for every root}

Consider a monic polynomial
\[
 H(Y)=Y^n+\sum_{j<n}(a_j+ib_j)Y^j
\]
over the represented field. Choose monic rational annihilators $p_j,q_j$
for its real and imaginary coefficient parts. Form the finite rational
algebra
\[
 T=\Q[U_0,V_0,\ldots,U_{n-1},V_{n-1},J]/
       \bigl(p_j(U_j),q_j(V_j),J^2+1\bigr).
\]
Its explicit monomial basis has dimension
$D_0=2\prod_{j<n}(\deg p_j)(\deg q_j)$.
The coefficient tuple defines an evaluation homomorphism sending
$U_j,V_j,J$ to $a_j,b_j,i$. No primitive-element theorem or compatible
embedding search is involved.

Let $M(Y)$ be the matrix of multiplication on this basis by
\[
 H_T(Y)=Y^n+\sum_{j<n}(U_j+JV_j)Y^j,
 \qquad B_0(Y)=\det M(Y).
\]
Since $M(Y)=Y^nI$ plus lower-degree terms, $B_0$ is monic of degree $nD_0$.
In particular it is nonzero, even though $T$ may have zero divisors.
Apply the adjugate matrix to the coordinate vector of $1$ to obtain an
element $E(Y)\in T[Y]$ satisfying
\[
 H_T(Y)E(Y)=B_0(Y)\,1.
\]
Evaluation at any named root $\beta$ of $H$ proves $B_0(\beta)=0$.
The construction is a finite determinant calculation with rational
coefficients. A nonmonic polynomial is first normalized by its certified
nonzero leading coefficient.

\subsection{Finite candidates and complete splitting}

Write $m=\deg B_0$. Define
\begin{align*}
 R_0(U)&=\Res_Z\bigl(B_0(Z),B_0(2U-Z)\bigr),\\
 T_0(V)&=\Res_Z\bigl(B_0(Z),B_0(Z-2iV)\bigr)\in\Q(i)[V],\\
 S_0(V)&=T_0(V)\overline{T_0}(V)\in\Q[V],
\end{align*}
where the bar conjugates coefficients.
The leading coefficients of $R_0,T_0$ are respectively
$2^{m^2}$ and $(-2i)^{m^2}$, and both have degree $m^2$.
Consequently $R_0$ and $S_0$ are nonzero.
These assertions follow by selecting the top parameter-degree terms in
the resultant determinant.

If $\beta=u+iv$ is a root of $B_0$, its conjugate is also a root because
$B_0$ has rational coefficients. Substitution at $Z=\beta$ in the
resultant identities gives
\[
 R_0(u)=0,\qquad T_0(v)=0,\qquad S_0(v)=0.
\]
Isolate all real roots $r_1,\ldots,r_A$ of $R_0$ and
$s_1,\ldots,s_B$ of $S_0$. The finite Cartesian list
\[
 \{\,r_j+i s_k:1\le j\le A,\ 1\le k\le B\,\}
\]
therefore contains every named complex root of $H$.
Completeness here includes a constructive finite identification of both
coordinates, by Lemma~\ref{lem:backend-isolation}.
Extraneous candidates caused by other algebraic embeddings are harmless.

Start with residual polynomial $H$. Process the candidate list in order.
At a candidate $\beta$, use exact equality to test the residual value.
Whenever it is zero, divide exactly by $Y-\beta$, record the root,
and repeat, for at most $n$ divisions in total.
This produces the exact identity
\[
 H(Y)=\prod_{\text{recorded }\beta}(Y-\beta)\,H_{\rm rem}(Y).
\]
The residual is monic. No processed candidate becomes a root again:
if it vanished after a later division, the preceding polynomial would
also have vanished there. Thus no candidate is a root of $H_{\rm rem}$.

Its degree is a decidable natural number. If that degree were positive,
the single-root constructive FTA of Section~\ref{sec:fta}, applied to
$H_{\rm rem}$, would supply a named root $\gamma$.
Its hypotheses hold: a single fixed primitive recursive evaluator reads
the algebraic coefficient codes, and the leading coefficient is $1$.
The displayed factorization implies $H(\gamma)=0$, so the candidate
construction identifies $\gamma$ with an already processed candidate.
This contradicts the residual tests. Hence $H_{\rm rem}=1$.
The recorded list is a complete splitting with multiplicities.

All executed loops are bounded: Euclidean steps by degree, the rational
grid by $2RK+1$, bisection by its explicit precision, determinants by their
finite dimensions, and synthetic divisions by $n$.
Finite code validity and the loop invariants are rational Boolean tests.
The correctness proof is therefore available to quantifier-free induction
after a precision horizon is fixed. Infinite names abbreviate their
explicit finite estimates. The reused FTA theorem supplies completeness
of this finite algorithm; it does not replace the exact coefficient
backend, and it does not introduce a universal evaluator for arbitrary
primitive recursive programs.


\section{The reused FTA result and source accounting}
\label{sec:fta}
The pinned \repo\ archive \texttt{docs/drafts/FTA\_iPRA.zip} supplies a
human-readable one-root theorem with this specialization:

\begin{quote}
For one fixed primitive recursive program giving coherent fast
Gaussian-rational coefficient names, a degree bound and positive evidence
that some nonconstant coefficient is apart from zero suffice to construct
a primitive recursive root name, with explicit Cauchy and residual moduli.
\end{quote}

For monic algebraic input the nonconstant coefficient witness is \(1\).
The program is the single fixed bisection-and-code evaluator constructed
in Section~\ref{sec:backend}. Polynomial codes and degrees are numerical
parameters, precisely as required by the dependency's schema.

That proof uses bounded rational descent: a direction routine from
square-root bisection and finite binomial identities; a finite comparison
of Taylor scales producing residual decrease; finite product estimates for
total travel; and an explicit sampling schedule for Cauchy and residual
moduli. The complete proof is frozen unchanged under
\texttt{dependencies/FTA\_iPRA/}, including its \texttt{.tex} and compiled
\texttt{.pdf}. The central one-root estimates were reviewed for this
application; \texttt{evidence/fta\_dependency\_audit.md} records that review.
This is mathematical review, not kernel verification.

The dependency's repeated numerical factorization discussion is only for
fixed degree and is not used as a uniform splitting algorithm.
Section~\ref{sec:backend} supplies the uniform exact construction.
The FTA draft's Lean source is an uncompiled specification; formalizing its
deep object derivation is an explicit remaining Lean dependency.

\subsection*{Provenance}
The repository sources were read at the following pinned commit:
\begin{center}\pin\end{center}
The precise record is \texttt{SOURCE\_MANIFEST.md}.
The FTA archive's Git blob identifier is
\texttt{60f15e1faa33b18ebbcb0a61f7fcccaa8b36aaf5}, and its size is
\(413411\) bytes.

The main report gives its own bounded algebraic reductions and analytic
estimates. The references identify the requested result, its constructive
context, and the reused foundation. They do not replace the finite
precision and logical-strength arguments.

\begin{thebibliography}{9}
\bibitem{mathworld}
E. W. Weisstein, \emph{Puiseux's Theorem}, MathWorld.
\url{https://mathworld.wolfram.com/PuiseuxsTheorem.html}.
Accessed 15 November 2026.
\bibitem{mannaa}
B. Mannaa and T. Coquand, Dynamic Newton--Puiseux theorem,
\emph{Journal of Logic and Analysis} \textbf{5}:5 (2013), 1--22.
\url{https://doi.org/10.4115/jla.2013.5.5}.
Open text: \url{https://math.hawaii.edu/home/jla/152-672-1-PB.pdf}.
\bibitem{mcmullen}
C. T. McMullen, \emph{Riemann Surfaces}, Harvard Math 213b course notes,
28 April 2026, Puiseux discussion, pp.~44--45.
\url{https://abel.math.harvard.edu/~ctm/math213b/home/course/course.pdf}.
Used for context and formulation, not as a constructive proof dependency.
\bibitem{proveit}
V. Reshetnikov and contributors, \emph{ProveIt}, snapshot
\texttt{d851229c33ef} (full commit in the source manifest).
Relevant paths:
\path{Logic/PrimitiveRecursiveArithmetic/ConciseIPRA/} and
\path{docs/drafts/FTA_iPRA.zip}.
\url{https://github.com/VladimirReshetnikov/ProveIt}.
The FTA source is a research draft with a human-readable proof and an
uncompiled Lean specification.
\end{thebibliography}
\end{document}
